\section{Related Work}
\textbf{Retrieval-Augmented Generation.}
RAG~\cite{DBLP:conf/nips/LewisPPPKGKLYR020} augments LLMs by injecting external
documents into prompts via similarity search, where memory is an unstructured
vector store and retrieval reduces to a one-shot query-time top-$k$ selection.
Two variants extend this paradigm. GraphRAG~\citep{graphrag} organizes the
retrieval corpus into graph structures and retrieves through community summaries
and neighborhood expansion, improving global and multi-hop reasoning over flat
similarity search. Agentic RAG instead moves retrieval inside the reasoning
loop. Search-o1~\cite{li2025searcho1} couples large reasoning models with an
agentic mechanism that issues search queries on demand whenever a knowledge gap
is detected and refines retrieved documents before integrating them into the
reasoning chain, and Search-R1~\cite{jin2025searchr1} trains this behavior
through reinforcement learning over multi-turn query generation. These methods
retrieve from open or external corpora to fill factual knowledge gaps within a
single reasoning task, rather than reconstructing over an agent's own persistent
interaction history.

\textbf{Graph-based Memory.}
Graph-based memory systems organize agent memory as a graph to capture
structural dependencies among memory units. A-Mem~\cite{DBLP:journals/corr/abs-2502-12110}
constructs structured memory notes linked through LLM-assisted relation
extraction and retrieves via seed selection followed by neighborhood expansion.
Zep~\cite{DBLP:journals/corr/abs-2501-13956} maintains a bi-temporal knowledge
graph that tracks when facts hold and invalidates outdated edges, supporting
retrieval over evolving knowledge. LiCoMemory~\cite{DBLP:journals/corr/abs-2511-01448}
organizes memory as a lightweight hierarchical graph that uses entities and
relations as a semantic indexing layer, with temporal and hierarchy-aware search
for efficient retrieval. While these representations improve relational and
multi-hop access, traversal is governed by predefined operators and does not
adapt to evidence accumulated during retrieval.

\textbf{Hierarchical and Persistent Memory Systems.}
Hierarchical memory systems maintain long-lived memory that is continually
updated across interactions. MemoryOS~\cite{DBLP:journals/corr/abs-2506-06326}
organizes memory into a short, mid, and long-term hierarchy for structured
access. Mem0~\cite{DBLP:journals/corr/abs-2504-19413} maintains a compact set of
salient facts through LLM-driven add, update, and delete operations.
SeCom~\cite{pan2025secom} constructs the memory bank at the level of topically coherent segments for more accurate retrieval.  These systems advance how
memory is stored and updated over time, but their retrieval procedures remain
passive, selecting memory units as a fixed function of the query without
reasoning over intermediate evidence.
% \textbf{Memory for LLM Agents}
% Retrieval-Augmented Generation (RAG)~\cite{DBLP:conf/nips/LewisPPPKGKLYR020} systems augment LLMs with external corpora by retrieving relevant documents via similarity search and injecting them into the prompt at inference time, forming the foundation of many memory-augmented agents. Building on this paradigm, GraphRAG \citep{graphrag} methods introduce explicit document- or entity-level graphs and retrieve information by expanding predefined neighborhoods from similarity-selected seeds. Extending RAG to long-horizon interactions, LLM agent memory systems model conversational experience as persistent memory with additional structure, such as graphs or hierarchical stores and support continual updates over time~\cite{DBLP:journals/corr/abs-2511-01448}. 
% Retrieval-Augmented Generation (RAG) \cite{DBLP:conf/nips/LewisPPPKGKLYR020} enhances LLMs by injecting external documents into prompts via similarity search. To handle complex dependencies, GraphRAG \citep{graphrag} and A-Mem \cite{DBLP:journals/corr/abs-2502-12110} organize information into graph structures, retrieving data through seed-based neighborhood expansion. Beyond static corpora, recent agent memory systems support continual updates \cite{DBLP:journals/corr/abs-2511-01448}: MemoryOS \cite{DBLP:journals/corr/abs-2506-06326} employs a hierarchical (short/mid/long-term) architecture for structured access, while Mem0 \cite{DBLP:journals/corr/abs-2504-19413} utilizes LLM-driven CRUD operations to maintain a compact set of salient facts for adaptive response generation

% A-Mem~\cite{DBLP:journals/corr/abs-2502-12110} represents LLM-generated memory notes as nodes in a directed memory graph and retrieves relevant notes through similarity-based seeding followed by neighborhood expansion.
% MemoryOS~\cite{DBLP:journals/corr/abs-2506-06326} organizes dialogue memory into a hierarchical structure consisting of short-term, mid-term, and long-term memory, and performs retrieval through hierarchical access.
% Mem0~\cite{DBLP:journals/corr/abs-2504-19413} maintains a compact set of salient facts through LLM-driven insertion, update, and deletion operations, and retrieves relevant facts using similarity-based search to condition response generation.